\chapter*{RINGKASAN}
\addcontentsline{toc}{chapter}{RINGKASAN}

Pengujian aplikasi web di Jurusan Teknologi Informasi Politeknik Negeri Malang menghadapi tantangan dalam memastikan kualitas aplikasi akademik seperti Kompen, Jayanti, dan Whistleblower. Alat automation testing seperti Selenium, Cypress, dan Playwright masih memerlukan pembuatan test case manual yang menjadi bottleneck. TathyaTest hadir sebagai test case generator otomatis yang mengintegrasikan web crawling (Colly), ekstraksi HTML (Goquery), dan dataset internal aplikasi untuk menghasilkan test case relevan dan kontekstual. Penelitian ini mengukur efektivitas TathyaTest dalam menghasilkan test case otomatis, meningkatkan cakupan pengujian fungsional (CRUD, autentikasi, akses modul), dan validasi ekspektasi berdasarkan dataset internal. Metode Prototyping digunakan secara iteratif dari prototipe CRUD sederhana, perbaikan berdasarkan umpan balik, hingga integrasi Playwright untuk cross-browser testing (Chrome, Firefox, WebKit) dengan dataset aplikasi internal. Penelitian dibatasi pada aplikasi web MVC dengan fokus pengujian fungsional berbasis UI. Analisis data menggunakan pendekatan kuantitatif (test coverage, efisiensi waktu, akurasi deteksi bug). Luaran meliputi prototipe TathyaTest TKT 4, publikasi jurnal Sinta 3, dan Hak Cipta. Penelitian ini berkontribusi pada Software Testing dan Quality Assurance melalui metode otomatis test case generation berbasis web crawling dan ekstraksi HTML, menjadi pondasi roadmap lima tahun pengembangan TathyaTest dengan integrasi AI dan machine learning.

\vspace{0.5cm}

\noindent
\textbf{Kata Kunci:} automation testing, test case generation, web crawling, functional testing, Playwright 
